\scanpage{85}
\mspage{84}
(iii) \deleted{Il n'existe pas de} \inserted{Pour \unsure{toute|tout} $\overline{F'}$}, avec
\[
F'\leqslant\overline{F'}\ll F
\]
on a $F'=\overline{F'}$, i.e $F'$ est maximal parmi les raffinements de $F$, ordonnés par $\leqslant$

\deleted{(iv) Pour tout $F'$ avec}
\[
\deleted{F'\ll\overline{F'}\ll F}
\]
\deleted{on a $F'=\overline{F'}$, i.e. $F'$ est maximal parmi les raffinements de $F$, ordonnés par $\ll$.}

\uline{Dém} \quad $\mathrm{(i)}\Longleftrightarrow\mathrm{(ii)}$ n'est autre que la prop. (\unsure{validée} par At 13), compte tenu de la remarque qui la suit.

\deleted{$\unsure{\mathrm{(ii)}|\mathrm{(iii)}}\Longrightarrow\mathrm{(iii)}$ est évident}

$\mathrm{(iii)}\Longrightarrow\mathrm{(ii)}$ \deleted{Car si \deleted{\unsure{on}} (ii) n'était \unsure{vraie|véri}} \inserted{i.e. $\overline{\mathrm{(ii)}}\Longrightarrow\overline{\mathrm{(iii)}}$.} En effet si $\deleted{\unsure{\exists x\in}}$ on a $x\in\operatorname{omb} F$, avec $\unsure{a|x}\parallel F$, cons\deleted{\unsure{é|i}}\unsure{idérons}
\[
F'\vee x=\overline{F'}
\]
prouvons que l'on a $\overline{F'}\ll F$ (ce qui montrera $\overline{\mathrm{(iii)}}$, car $F'\unsure{\lneq|\not\leqslant|\neq}\overline{F'}$). Or \deleted{\unsure{il suffit|il convient} donc de prouver cela} par At 4 cela se prouve strate par strate, on gagne.

\deleted{\uline{Prop} \unsure{Soient|Soit} $F$ figure,}

Il reste à prouver que $\mathrm{(ii)}\Longrightarrow\unsure{\mathrm{(iv)}|\mathrm{(iii)}}$. \deleted{Prouvons d'abord $\mathrm{(ii)}\Longrightarrow\mathrm{(iii)}$.} Si dans (iii), on avait $F'\neq\overline{F'}$, soit $\deleted{\unsure{x}}X\in\widetilde{\overline{F'}}\setminus\widetilde{F'}$, $x\in\operatorname{omb}(X)\unsure{\}|y}$ alors $x\ll\overline{F'}\ll F$ donc $x\ll F$, et $x\parallel F'$, \unsure{contrairemt|contrairement} à (ii). \deleted{Ainsi (i)(ii)(iii) sont \unsure{équivalents|équivalentes} et impliqués par (iv).} Mais impliquent ils (iv) ?
